\documentclass[11pt]{article}
\usepackage[utf8]{inputenc}
\usepackage{amsmath, amssymb, amsthm}
\usepackage[margin=1in]{geometry}

\theoremstyle{definition}
\newtheorem{definition}{Definition}
\newtheorem{theorem}{Theorem}
\newtheorem{lemma}{Lemma}
\newtheorem{proposition}{Proposition}

\theoremstyle{remark}
\newtheorem{remark}{Remark}

\DeclareMathOperator{\Id}{Id}

\title{Lecture 6: Continuous Modifications, Additive Noise, and Weak Topologies}
\date{20.11.2025}
\author{Tien Anh Vu}

\begin{document}
\maketitle

This lecture begins with the problem of path continuity for stochastic convolutions, proves the existence of a continuous modification by the factorization method, and then turns to additive noise, weak topologies, and compactness principles that are central to the analysis of stochastic evolution equations.

\section*{1. Continuous Modifications of the Stochastic Convolution}

\subsection*{1.1. The SPDE and the Non-Continuous Convolution}
We begin by restating our standard Stochastic Partial Differential Equation evolving in our Hilbert space \(H\):
\[
dX_t=(AX_t+B(s,X_s))\,dt+C(s,X_s)\,dW_t,
\qquad
X_0=\xi.
\]
Here the drift and diffusion coefficients explicitly depend on time through \(s\).

Assuming the linear operator \(A\) generates a \(C_0\)-semigroup \((T_t)_{t\ge 0}\), we immediately write down its corresponding mild solution:
\begin{align*}
X_t
&=
T_t\xi
+\underbrace{\int_0^t T_{t-s}B(s,X_s)\,ds}_{\text{drift}}
+\underbrace{\int_0^t T_{t-s}C(s,X_s)\,dW_s}_{\text{non continuous convolution}}.
\end{align*}
The first integral is the drift term. It is deterministic once the path \(X_s\) is fixed, and it collects the accumulated effect of the nonlinear drift coefficient \(B(s,X_s)\) propagated by the semigroup.

The second integral is the stochastic convolution. Because the semigroup operator \(T_{t-s}\) depends on the evaluation time \(t\) and remains inside the stochastic integral, continuity is not immediate. Resolving this is the main objective.

\subsection*{1.2. Proposition 2.2: Existence of a Continuous Modification}
To resolve this difficulty, we now turn to Proposition 2.2 on the existence of a continuous modification.

\begin{proposition}
Let \(p>1\), and let our integrand \(\Phi\) belong to the space
\[
\Phi\in L^{2p}(\Omega_T,\mathcal P_T,P_T;L_2(U,H)).
\]
Assume moreover that
\[
\mathbb E\!\left(\int_0^T \|\Phi(s)\circ Q^{1/2}\|_{L_2(U,H)}^{2p}\,ds\right)<\infty.
\]
Then there exists a constant \(C_T\) such that
\begin{align*}
\mathbb E\!\left(
\sup_{t\in[0,T]}
\left\|
\int_0^t T_{t-s}\Phi(s)\,dW_s
\right\|_H^{2p}
\right)
&\le
C_T\,
\mathbb E\!\left(\int_0^T \|\Phi(s)\circ Q^{1/2}\|_{L_2(U,H)}^{2p}\,ds\right),
\end{align*}
and, moreover,
\[
W_A^\Phi(t):=\int_0^t T_{t-s}\Phi(s)\,dW_s
\]
has a continuous modification.
\end{proposition}

Here \(\Phi\in L^{2p}(\Omega_T,\mathcal P_T,P_T;L_2(U,H))\) means that \(\Phi\) is progressively measurable on \([0,T]\times\Omega\), takes values in the Hilbert-Schmidt operators from \(U\) to \(H\), and has the required \(2p\)-integrability in time and probability.

This extends the maximal inequalities studied previously by giving both a \(2p\)-moment estimate and a continuous modification.

\subsection*{1.3. Proof Strategy: The Factorization Method}
To prove this proposition, we use the factorization method.

To mathematically factorize the semigroup out of the convolution, we must utilize a specific fractional integral identity. For any times \(s<t\), we evaluate the following integral:
\begin{align*}
\int_s^t (t-r)^{\alpha-1}(r-s)^{-\alpha}\,dr
&=
\frac{1}{C_\alpha}
&=
\frac{\Gamma(1)}{\Gamma(\alpha)\Gamma(1-\alpha)}
&=
\frac{\sin(\pi\alpha)}{\pi}.
\end{align*}
The second and third equalities rewrite the constant using the Gamma function and Euler's reflection formula.

By strategically inserting this constant identity into the stochastic integral, we will be able to split the time domain, extract the semigroup, and prove the continuous modification of our SPDE solution.

\section*{2. The Factorization Method in Detail}

We now prove Proposition 2.2. More precisely, we must show that the stochastic convolution \(W_{A,\Phi}\) satisfies the required \(2p\)-moment estimate and admits a continuous modification.

\begin{proof}
\noindent\textbf{2.1. Factorization and Stochastic Fubini.}
We begin by confirming the substitution
\[
u:=\frac{r-s}{t-s}.
\]
This maps the Beta integral to the unit interval
\[
\int_0^1 (1-u)^{\alpha-1}u^{-\alpha}\,du=C_\alpha,
\]
which gives the constant \(C_\alpha\).

With this constant, we rewrite the stochastic convolution as
\begin{align*}
W_{A,\Phi}(t)
&=
\int_0^t T_{t-s}(\Phi(s))\,dW_s \\
&=
C_\alpha\int_0^t\int_s^t (t-r)^{\alpha-1}(r-s)^{-\alpha}\,dr\,T_{t-s}(\Phi(s))\,dW_s.
\end{align*}
Next, we deploy a crucial property of our \(C_0\)-semigroup. We split the temporal operator according to
\[
T_{t-s}=T_{t-r}T_{r-s}.
\]
This separates the time variables into \((t-r)\) and \((r-s)\).

By swapping the order of integration, we obtain
\begin{align*}
W_{A,\Phi}(t)
&=
C_\alpha\int_0^t\int_0^r (t-r)^{\alpha-1}(r-s)^{-\alpha}T_{t-r}T_{r-s}\Phi(s)\,dW_s\,dr \\
&=
C_\alpha\int_0^t (t-r)^{\alpha-1}T_{t-r}\left(\int_0^r (r-s)^{-\alpha}T_{r-s}\Phi(s)\,dW_s\right)\,dr.
\end{align*}
Define the auxiliary process by
\begin{align*}
y(r)
&:=
\int_0^r (r-s)^{-\alpha}T_{r-s}\Phi(s)\,dW_s.
\end{align*}
Then
\begin{align*}
W_{A,\Phi}(t)
&=
C_\alpha\int_0^t (t-r)^{\alpha-1}T_{t-r}y(r,\omega)\,dr.
\end{align*}
This is a standard Lebesgue integral in time, so the resulting process is continuous in \(t\).

\medskip

\noindent\textbf{2.2. The \(2p\)-Moment Bound.}
To make this rigorous, we also need a bound on the expected norm. Assume that
\[
p>\frac{1}{2\alpha}>1.
\]
Applying the Bochner inequality and then H\"older's inequality yields
\begin{align*}
\|W_{A,\Phi}(t)\|_H^{2p}
&=
\left\|C_\alpha\int_0^t (t-r)^{\alpha-1}T_{t-r}y(r,\omega)\,dr\right\|_H^{2p} \\
&\le
C_\alpha^{2p}\left(\int_0^t (t-r)^{\alpha-1}\|T_{t-r}(y(r))\|_H\,dr\right)^{2p} \\
&\le
C_\alpha^{2p}
\left(\int_0^t (t-r)^{(\alpha-1)\frac{2p}{2p-1}}\,dr\right)^{2p-1}
\cdot \int_0^t \|T_{t-r}(y(r))\|_H^{2p}\,dr.
\end{align*}

The first integral is computed by
\begin{align*}
\int_0^t (t-r)^{(\alpha-1)\frac{2p}{2p-1}}\,dr
&=
\frac{1}{1+(\alpha-1)\frac{2p}{2p-1}}\,t^{\,1+(\alpha-1)\frac{2p}{2p-1}}.
\end{align*}
Since \(p>\frac{1}{2\alpha}\), the exponent \((\alpha-1)\frac{2p}{2p-1}\) is greater than \(-1\), so this term is bounded by a finite constant depending only on \(T\).

Absorbing all constants into \(C_{\alpha,p,T}\), we obtain
\begin{align*}
\|W_{A,\Phi}(t)\|_H^{2p}
&\le
C_\alpha^{2p}
\left(\int_0^t (t-r)^{(\alpha-1)\frac{2p}{2p-1}}\,dr\right)^{2p-1}
\int_0^t \|T_{t-r}(y(r))\|_H^{2p}\,dr \\
&\le
C_{\alpha,p,T}\int_0^t \|T_{t-r}y(r)\|_H^{2p}\,dr.
\end{align*}

\medskip

\noindent\textbf{3.1. Integrability of the Auxiliary Process.}
It remains to check the integrability of the final term. Since
\[
y(r)=\int_0^r (r-s)^{-\alpha}T_{r-s}\Phi(s)\,dW_s,
\]
the Burkholder-Davis-Gundy inequality gives
\begin{align*}
\mathbb E(\|y(r)\|_H^{2p})
&=
\mathbb E\!\left(\left\|\int_0^r (r-s)^{-\alpha}T_{r-s}\Phi(s)\,dW_s\right\|_H^{2p}\right) \\
&\le
C_p\,\mathbb E\!\left(\left(\int_0^r (r-s)^{-2\alpha}\|\Phi(s)\circ Q^{1/2}\|_{L_2(U,H)}^2\,ds\right)^p\right).
\end{align*}
Integrating over \(r\in[0,T]\), swapping expectation and the time integral, and then applying H\"older's inequality to the inner integral
\[
\int_0^r (r-s)^{-2\alpha}\|\Phi(s)\circ Q^{1/2}\|_{L_2(U,H)}^2\,ds
\]
yields
\begin{align*}
\int_0^T \mathbb E(\|y(r)\|_H^{2p})\,dr
&\le
C_p\,\mathbb E\!\left(\int_0^T \left(\int_0^r (r-s)^{-2\alpha}\|\Phi(s)\circ Q^{1/2}\|_{L_2(U,H)}^2\,ds\right)^p\,dr\right) \\
&\le
C_p\,\mathbb E\!\left[\int_0^T \left(\int_0^r (r-s)^{-2\alpha}\,ds\right)^{p-1}
\int_0^r (r-s)^{-2\alpha}\|\Phi(s)\circ Q^{1/2}\|_{L_2(U,H)}^{2p}\,ds\,dr\right].
\end{align*}
Since \(\alpha<\frac12\),
\begin{align*}
\int_0^r (r-s)^{-2\alpha}\,ds
&\le
\frac{1}{1-2\alpha}T^{1-2\alpha},
\end{align*}
so all constants can be absorbed into \(\widetilde C_{p,\alpha,T}\), and therefore
\begin{align*}
\int_0^T \mathbb E(\|y(r)\|_H^{2p})\,dr
&\le
\widetilde C_{p,\alpha,T}\,\mathbb E\!\left(\int_0^T \|\Phi(s)\circ Q^{1/2}\|_{L_2(U,H)}^{2p}\,ds\right)
<\infty.
\end{align*}
Hence \(y\in L^{2p}([0,T];H)\).
Recall from \(2.2\) that
\[
\|W_{A,\Phi}(t)\|_H^{2p}
\le
C_{\alpha,p,T}\int_0^t \|T_{t-r}y(r)\|_H^{2p}\,dr.
\]
Using the boundedness of the semigroup on \([0,T]\), we obtain
\[
\|W_{A,\Phi}(t)\|_H^{2p}
\le
C_{\alpha,p,T}\sup_{0\le u\le t}\|T_u\|_{L(H)}^{2p}\int_0^t \|y(r)\|_H^{2p}\,dr.
\]
Therefore
\[
\sup_{t\in[0,T]}\|W_{A,\Phi}(t)\|_H^{2p}
\le
C_{\alpha,p,T}\sup_{0\le u\le T}\|T_u\|_{L(H)}^{2p}\int_0^T \|y(r)\|_H^{2p}\,dr.
\]
Taking expectation, we obtain
\begin{align*}
\mathbb E\!\left(
\sup_{t\in[0,T]}
\|W_{A,\Phi}(t)\|_H^{2p}
\right)
&\le
C_{\alpha,p,T}\sup_{0\le u\le T}\|T_u\|_{L(H)}^{2p}\int_0^T \mathbb E(\|y(r)\|_H^{2p})\,dr \\
&\le
C_{\alpha,p,T}\sup_{0\le u\le T}\|T_u\|_{L(H)}^{2p}\widetilde C_{p,\alpha,T}
\mathbb E\!\left(\int_0^T \|\Phi(s)\circ Q^{1/2}\|_{L_2(U,H)}^{2p}\,ds\right) \\
&\le
C_T\,\mathbb E\!\left(\int_0^T \|\Phi(s)\circ Q^{1/2}\|_{L_2(U,H)}^{2p}\,ds\right).
\end{align*}
This is exactly the required maximal estimate.

\medskip

\noindent\textbf{3.2. Main Implication.}
Define
\[
F(y)(t):=\int_0^t (t-r)^{\alpha-1}T_{t-r}y(r)\,dr.
\]
The factorization formula shows that
\[
W_{A,\Phi}(t)=C_\alpha F(y)(t).
\]
Since \(y\in L^{2p}([0,T];H)\) and the kernel \((t-r)^{\alpha-1}\) is integrable, the map \(F\) sends \(L^{2p}([0,T];H)\) into \(C([0,T],H)\). Therefore the stochastic convolution \(W_{A,\Phi}\) admits a continuous modification.

\noindent This proves that the stochastic convolution \(W_{A,\Phi}\) admits a continuous modification and satisfies the required \(2p\)-moment bound.
\end{proof}

\section*{4. Additive Noise and Weak Topology}

Having established the existence of a continuous modification for the stochastic convolution, we now turn to SPDEs with additive noise and then to the weak topological framework needed for stability.

\subsection*{4.1. SPDEs with Additive Noise: The Pathwise Decomposition}
We introduce the special case of stochastic differential equations with additive noise. Our governing equation is
\[
dX_t=[AX_t+B(X_t)]\,dt+C\,dW_t.
\]
The diffusion operator \(C\) does not depend on the state \(X_t\). This is exactly the additive-noise case.

Because the noise is additive, we can isolate the purely stochastic part by defining the stochastic convolution
\[
W_{A,C}(t):=\int_0^t T_{t-s}C\,dW_s.
\]
We then decompose the entire state process into two distinct parts:
\[
X_t=Y_t+W_{A,C}(t).
\]
By substituting this decomposition back into our original SPDE, the stochastic differential \(C\,dW_t\) cancels out, leaving us with a new equation for \(Y_t\):
\[
dY_t=AY_t\,dt+B(Y_t+W_{A,C}(t))\,dt,
\qquad
Y_0=\xi.
\]
This removes the stochastic differential from the equation for \(Y_t\). The randomness is absorbed into the drift term as a known pathwise input.

The mild solution of this transformed system is
\[
Y_t=T_t\xi+\int_0^t T_{t-s}B(Y_s+W_{A,C}(s))\,ds.
\]
Adding the stochastic convolution back in, we recover the original state:
\[
X_t=T_t\xi+\int_0^t T_{t-s}B(Y_s+W_{A,C}(s))\,ds+W_{A,C}(t).
\]
This pathwise formulation is often analytically simpler.

\subsection*{4.2. Interlude: Appendix B and the Weak Topology}
To prepare for the upcoming proofs regarding the stability and convergence of these solutions, we now turn to an interlude on weak topology.

Let \(E\) be a real Banach space. To understand the weak topology, we must simultaneously consider its continuous dual space \(E'\), which consists of all bounded linear functionals acting on \(E\).

We define weak convergence as follows: a sequence \((x_n)_{n\in\mathbb N}\subset E\) converges weakly to \(x\in E\) if, for every functional \(f\in E'\),
\[
\lim_{n\to\infty}\langle f,x_n\rangle=\langle f,x\rangle.
\]
Here \(\langle f,x\rangle\) denotes the dual pairing, that is, the action of the functional \(f\) on the vector \(x\). In this case we write
\[
x_n\rightharpoonup x.
\]

\subsection*{4.3. Connecting Strong and Weak Convergence}
We finally establish the basic hierarchy of convergence: strong convergence with respect to the norm implies weak convergence.

Suppose that
\[
\|x_n-x\|_E\to 0.
\]
We must show that
\[
\langle f,x_n\rangle\to\langle f,x\rangle
\qquad \text{for every } f\in E'.
\]
By linearity of \(f\),
\[
|\langle f,x_n\rangle-\langle f,x\rangle|
=
|\langle f,x_n-x\rangle|.
\]
Using the fundamental dual estimate, we obtain
\[
|\langle f,x_n-x\rangle|
\le
\|f\|_{E'}\|x_n-x\|_E.
\]
Because \(\|f\|_{E'}\) is a finite constant and \(\|x_n-x\|_E\to 0\), the right-hand side converges to \(0\). Therefore
\[
x_n\rightharpoonup x.
\]
Thus strong convergence implies weak convergence.

\section*{5. Weak Topologies Beyond the First Implication}

We continue the discussion of weak convergence by examining why the converse fails, why weak limits are unique, and how strong convergence can be recovered in Hilbert spaces when the norms also converge.

\subsection*{5.1. The Converse is Not True}
We now address the converse question: does weak convergence imply strong convergence? The answer is no.

To prove this, we examine a classic counterexample in the sequence space \(\ell^2\). Consider the standard orthonormal basis sequence
\[
e_n=(\delta_{kn})_{k\in\mathbb N}.
\]
Because each vector is a unit basis vector,
\[
\|e_n\|_{\ell^2}=1
\qquad \text{for all } n,
\]
so the sequence does not converge strongly to \(0\).

However, the sequence does converge weakly to \(0\). For any functional \(f\in (\ell^2)'\), the Riesz representation theorem identifies \(f\) with a sequence \((f_n)_{n\in\mathbb N}\in\ell^2\), and
\[
\langle f,e_n\rangle_{\ell^2}=f_n.
\]
Since \(f\in\ell^2\), we have \(f_n\to 0\), and therefore
\[
e_n\rightharpoonup 0.
\]
Thus weak convergence does not imply strong convergence.

\subsection*{5.2. Uniqueness of the Weak Limit and Hahn-Banach}
We next record that the weak limit is unique.

If a sequence were to converge weakly to two different limits, their difference, say \(x\), would satisfy
\[
\langle f,x\rangle=0
\qquad \text{for all } f\in E'.
\]
This implies \(x=0\), hence the two limits must coincide. The key reason is the Hahn-Banach theorem, which guarantees that for every nonzero \(x\in E\), there exists a functional \(f\in E'\) such that
\[
f(x)\neq 0.
\]
By contraposition, if every functional vanishes on a vector, that vector must be zero.

The structural content of Hahn-Banach is the extension of bounded linear functionals. If \(l\) is a continuous linear functional defined on a subspace \(E_0\subset E\), then there exists an extension \(\overline l\in E'\) such that
\[
\overline l|_{E_0}=l
\qquad \text{and} \qquad
\|\overline l\|_{E'}=\|l\|.
\]
This extension principle allows one to separate points and guarantee uniqueness of weak limits.

\subsection*{5.3. Weak Lower Semicontinuity of the Norm}
We now turn to the behavior of norms under weak convergence. If
\[
x_n\rightharpoonup x
\qquad \text{in } E,
\]
then \((x_n)\) is bounded and
\[
\|x\|_E\le\liminf_{n\to\infty}\|x_n\|_E.
\]

For every \(f\in E'\), weak convergence gives
\[
|\langle f,x\rangle|
=
\lim_{n\to\infty}|\langle f,x_n\rangle|.
\]
Applying the standard dual bound yields
\[
|\langle f,x\rangle|
\le
\liminf_{n\to\infty}\|f\|_{E'}\|x_n\|_E.
\]
By Hahn-Banach, we may choose a functional \(f\in E'\) with
\[
\|f\|_{E'}=1
\qquad \text{and} \qquad
\langle f,x\rangle=\|x\|_E.
\]
Substituting this specific functional into the inequality above gives
\[
\|x\|_E\le\liminf_{n\to\infty}\|x_n\|_E.
\]

\subsection*{5.4. Recovering Strong Convergence}
Finally, we come to the standard Hilbert space criterion for recovering strong convergence: weak convergence together with convergence of the norms implies strong convergence.

Assume that
\[
x_n\rightharpoonup x
\qquad \text{and} \qquad
\|x_n\|\to\|x\|.
\]
We expand the squared distance:
\[
\|x_n-x\|^2
=
\|x_n\|^2-2\langle x,x_n\rangle+\|x\|^2.
\]
As \(n\to\infty\), the first term converges to \(\|x\|^2\) by assumption. Since \(y\mapsto \langle x,y\rangle\) is a continuous linear functional, weak convergence gives
\[
\langle x,x_n\rangle\to\langle x,x\rangle=\|x\|^2.
\]
Therefore
\[
\|x_n-x\|^2
\to
\|x\|^2-2\|x\|^2+\|x\|^2
=
0.
\]
Hence
\[
x_n\to x
\]
strongly.

This concludes the discussion of weak and strong convergence.

\section*{6. Compactness and Weak-* Convergence}

We now complete the discussion of topological convergence by introducing compactness in the weak-* topology and connecting it to probability measures through the Riesz-Radon theorem.

\subsection*{6.1. The Conclusion of Strong Convergence}
We begin with the definitive algebraic conclusion of the proof from the previous session:
\[
\|x\|^2-2\langle x,x\rangle+\|x\|^2=0.
\]
Since \(\langle x,x\rangle=\|x\|^2\), the three terms cancel, and this gives strong convergence.

\subsection*{6.2. The Banach-Alaoglu Theorem}
We now introduce the Banach-Alaoglu theorem, the infinite-dimensional analogue of Bolzano-Weierstrass. In infinite dimensions, compactness is recovered only in a weaker topology.

The theorem states that the unit ball, and hence any bounded subset, in \(E'\), where \(E'\) is the dual of a separable Banach space \(E\), is sequentially compact in the weak-* topology. Thus every bounded sequence of functionals \((f_n)\) has a weak-* convergent subsequence. This means that for every fixed vector \(x\in E\),
\[
\langle f_n,x\rangle\to \langle f,x\rangle.
\]
So bounded families of functionals admit weak-* limits.

\subsection*{6.3. Application to Reflexive and Hilbert Spaces}
We next turn to reflexive spaces. A space is reflexive if its bidual equals itself:
\[
E''=E.
\]
The most important examples of this, as noted in the lecture, are standard Hilbert spaces.

In reflexive spaces, weak-* compactness in the bidual becomes weak compactness in the space itself. Consequently, any bounded subset of \(E\) is sequentially weakly compact. In particular, every bounded sequence in a Hilbert space admits a weakly convergent subsequence.

\subsection*{6.4. Remark: Connecting to Probability Theory}
We now connect this framework to probability theory. Consider weak convergence of probability measures, denoted
\[
\mu_n \rightharpoonup \mu
\qquad \text{in probability}.
\]
This is defined through testing against continuous functions: if \(\Omega\subset \mathbb R^d\) is closed and bounded, then weak convergence means that for every \(f\in C(\Omega)\),
\[
\int f\,d\mu_n \to \int f\,d\mu.
\]
This can be written as
\[
\langle \mu_n,f\rangle \to \langle \mu,f\rangle,
\]
This is convergence in the functional analytic sense, because integration against a measure is a bounded linear functional on \(C(\Omega)\). Thus probabilistic weak convergence is weak-* convergence in the dual of \(C(\Omega)\).

\subsection*{6.5. The Riesz-Radon Theorem}
This connection is formalized by the Riesz-Radon theorem. It states that the continuous dual space of \(C(\Omega)\) is exactly the space of all finite signed measures:
\[
C(\Omega)'=\{\text{finite signed measures on }\Omega\}.
\]
This identifies probability measures with bounded linear functionals on continuous test functions. Weak convergence of measures can therefore be studied through weak-* compactness in the dual of \(C(\Omega)\).

This completes the discussion of compactness and weak-* convergence.

\section*{Appendix}

\subsection*{A.1. The Beta Function}
\begin{definition}
For parameters \(a,b>0\), the Beta function is defined by
\[
B(a,b):=\int_0^1 u^{a-1}(1-u)^{b-1}\,du
=
\frac{\Gamma(a)\Gamma(b)}{\Gamma(a+b)}.
\]
\end{definition}

\subsection*{A.2. The Stochastic Fubini Theorem}
\begin{definition}
The Stochastic Fubini theorem is the result that permits an interchange between a stochastic integral and an ordinary integral, provided the required measurability and integrability assumptions are satisfied. In the present setting, its role is to justify identities of the form
\[
\int_0^t\int_s^t f(r,s)\,dr\,dW_s
=
\int_0^t\int_0^r f(r,s)\,dW_s\,dr.
\]
\end{definition}

\subsection*{A.3. Continuous Modification}
\begin{definition}
A process \(Y\) is said to admit a continuous modification if there exists another process \(\widetilde Y\) with continuous sample paths such that
\[
P\bigl(Y_t=\widetilde Y_t \text{ for all } t\in[0,T]\bigr)=1.
\]
\end{definition}

\subsection*{A.4. Additive Noise}
\begin{definition}
An SPDE is said to have additive noise if the diffusion coefficient does not depend on the state variable. A typical form is
\[
dX_t=[AX_t+B(X_t)]\,dt+C\,dW_t,
\]
where \(C\) is a fixed operator.
\end{definition}

\subsection*{A.5. Weak Convergence}
\begin{definition}
Let \(E\) be a Banach space. A sequence \((x_n)\subset E\) converges weakly to \(x\in E\) if
\[
\langle f,x_n\rangle\to \langle f,x\rangle
\qquad \text{for every } f\in E'.
\]
In that case one writes
\[
x_n\rightharpoonup x.
\]
\end{definition}

\subsection*{A.6. Weak-* Convergence}
\begin{definition}
Let \(E\) be a Banach space and let \((f_n)\subset E'\). One says that \(f_n\) converges weakly-* to \(f\in E'\) if
\[
\langle f_n,x\rangle\to \langle f,x\rangle
\qquad \text{for every } x\in E.
\]
\end{definition}

\subsection*{A.7. Reflexive Spaces}
\begin{definition}
A Banach space \(E\) is called reflexive if the canonical embedding of \(E\) into its bidual \(E''\) is surjective, so that
\[
E''=E.
\]
Hilbert spaces are the main examples of reflexive spaces used in these notes.
\end{definition}

\subsection*{A.8. The Riesz-Radon Identification}
\begin{theorem}
If \(\Omega\subset \mathbb R^d\) is compact, then the Riesz-Radon theorem identifies the dual of \(C(\Omega)\) with the space of finite signed measures on \(\Omega\):
\[
C(\Omega)'=\{\text{finite signed measures on }\Omega\}.
\]
This is the functional-analytic basis of weak convergence of probability measures.
\end{theorem}

\subsection*{A.9. The Hahn-Banach Theorem}
\begin{theorem}
Let \(E\) be a normed space, let \(E_0\subset E\) be a linear subspace, and let \(l:E_0\to \mathbb R\) be a bounded linear functional. Then there exists an extension \(\overline l:E\to \mathbb R\) such that
\[
\overline l|_{E_0}=l
\qquad \text{and} \qquad
\|\overline l\|_{E'}=\|l\|_{E_0'}.
\]
\end{theorem}
This theorem is used to separate points by bounded linear functionals. In particular, it explains why a weak limit is unique and why one can choose a functional that attains the norm in the proof of lower semicontinuity.

\subsection*{A.10. The Banach-Alaoglu Theorem}
\begin{theorem}
Let \(E\) be a Banach space. The closed unit ball of the dual space \(E'\) is compact in the weak-* topology. In the separable setting used in these notes, this yields sequential weak-* compactness for bounded subsets of \(E'\).
\end{theorem}
This is the compactness principle used to extract weak-* convergent subsequences from bounded families of functionals. It plays the role of Bolzano-Weierstrass in infinite dimensions.

\subsection*{A.11. Weak Lower Semicontinuity of the Norm}
\begin{lemma}
If \(x_n\rightharpoonup x\) weakly in a Banach space \(E\), then
\[
\|x\|_E\le \liminf_{n\to\infty}\|x_n\|_E.
\]
\end{lemma}
This lemma shows that the norm cannot jump upward under weak convergence. It is one of the main tools for passing to limits in bounded sequences and for comparing weak and strong convergence.

\end{document}
